% Generated draft: requires mathematical and source review.
\documentclass[11pt]{article}
\usepackage[margin=25mm]{geometry}
\usepackage{fontspec}
\setmainfont{Latin Modern Roman}
\usepackage{amsmath,amssymb,mathtools,unicode-math}
\setmathfont{Latin Modern Math}
\usepackage{enumitem,needspace,xurl,hyperref}
\setlength{\parindent}{0pt}
\setlength{\parskip}{5pt}
\setlength{\emergencystretch}{3em}
\newcommand{\N}{\mathbb{N}}
\newcommand{\Z}{\mathbb{Z}}
\newcommand{\Q}{\mathbb{Q}}
\newcommand{\R}{\mathbb{R}}
\newcommand{\C}{\mathbb{C}}
\newcommand{\F}{\mathbb{F}}
\newcommand{\A}{\mathbb{A}}
\newcommand{\PP}{\mathbb{P}}
\newcommand{\E}{\mathbb{E}}
\newcommand{\eps}{\varepsilon}
\newcommand{\dd}{\,\mathrm{d}}
\begin{document}
\section*{UnsolvedMath ID: 801061}
\textbf{Source problem code:} catalogue entry
\subsection{Definitions and mathematical statement}
Let \(\mathcal{WY}\) denote the parameter locus named in the supplied source, and let \(H_{a,b}\) denote its H\'enon map for \((a,b)\in\mathcal{WY}\). Write \(\mathfrak K_{a,b}\) for the collection of kneading sequences denoted by \(\mathfrak K\) for this map in the source. The record does not define \(\mathcal{WY}\), specify the formula or normalization of \(H_{a,b}\), or define these kneading sequences and their symbolic coding. These remain unresolved prerequisites; in particular, no real or complex parameter domain is stipulated here.

\textbf{Question 13 (Berger--Bedford, 2023).} Under the source's intended definitions, does every pair of distinct kneading sequences belonging to a map with parameters in \(\mathcal{WY}\) determine its ordered parameter pair uniquely within \(\mathcal{WY}\)? If the sequences for different parameters are compared in a common symbolic coding, this asks whether
\[
\begin{gathered}
\text{for all }(a,b),(a',b')\in\mathcal{WY}
\text{ and all }s,t\in\mathfrak K_{a,b}\cap\mathfrak K_{a',b'},\\
s\ne t\quad\Longrightarrow\quad(a,b)=(a',b').
\end{gathered}
\]
Existence of a realizing parameter pair is assumed by taking the two sequences from one such map; the question does not assert realizability of arbitrary symbolic pairs. The source further asks whether the two sequences thereby govern every other kneading sequence in \(\mathfrak K_{a,b}\). This can be read as determination of the entire collection through uniqueness of the parameters, but the record does not specify whether ``govern'' requires any further reconstruction rule.
\subsection{Short English statement}
Berger and Bedford ask whether any two distinct kneading sequences of a H\'enon map with parameters in \(\mathcal{WY}\) uniquely identify those parameters and thereby determine all its other kneading sequences. The parameter locus, map normalization, symbolic coding, and meaning of ``govern'' require clarification from the source.
\subsection{Sources}
\begin{itemize}
\item Local UnsolvedMath registry: permanent ID \(801061\), source problem code \texttt{catalogue entry}, title \emph{Henon Maps: Open Problems --- Question 13}. The supplied registry attributes the problem to Berger and Bedford, with proposed year 2023, and labels it \texttt{open}; that label is supplied metadata, not an independently verified mathematical or literature assessment.
\item Berger and Bedford, \emph{Henon Maps: A List of Open Problems} (2023), Question 13. Supplied source URL: \url{https://arxiv.org/abs/2312.03907}. The source record reports access on 2026-10-16 and extraction by \texttt{arxiv\_latex}; these are supplied provenance statements, not access or verification performed for this draft.
\end{itemize}
\end{document}
